% Notes: loose structure, quick to write in. The preamble (up to
% \begin{document}) gives them their look.
\documentclass[11pt,a4paper]{article}

% Fonts and typography
\usepackage[T1]{fontenc}
\usepackage{lmodern}
\usepackage{microtype}

% Page: paragraphs separated by space, not indented
\usepackage[a4paper,hmargin=2.4cm,vmargin=2.6cm]{geometry}
\usepackage{parskip}

% Mathematics and figures
\usepackage{amsmath,amssymb}
\usepackage{graphicx}
\graphicspath{{figures/}}

% Citations, author-year (natbib with BibTeX)
\usepackage[round]{natbib}

% Links and citations in the accent colour
\usepackage{xcolor}
\definecolor{accent}{RGB}{140,82,18}
\usepackage[colorlinks=true,linkcolor=accent,citecolor=accent,urlcolor=accent]{hyperref}

% Unnumbered sections
\setcounter{secnumdepth}{0}

\newcommand{\notestitle}{{{title}}}
\newcommand{\notesauthor}{{{author}}}
\newcommand{\notesdate}{{{date}}}

\hypersetup{pdftitle={\notestitle},pdfauthor={\notesauthor}}

\begin{document}

{\LARGE\bfseries \notestitle\par}
\vspace{0.3em}
{\color{black!60}\notesauthor\quad\textperiodcentered\quad\notesdate\par}
\vspace{0.4em}
\noindent\rule{\linewidth}{0.4pt}

\section{Ideas}

\begin{itemize}
  \item Write things down as they come; structure can wait.
  \item One idea per bullet, so it is easy to move later.
  \begin{itemize}
    \item Sub-points go one level in.
  \end{itemize}
\end{itemize}

\section{Reading}

Keep track of what you read and what it gave you. Programs, like papers,
should be written for people first~\citep{knuth1984}. Every source cited from
\texttt{refs.bib} appears in the list at the end.

\section{Next steps}

\begin{enumerate}
  \item Turn the strongest idea into a paragraph.
  \item Ask the agent to find the related work and add it to
    \texttt{refs.bib}.
  \item Start a paper from these notes when the outline holds.
\end{enumerate}

\bibliographystyle{plainnat}
\bibliography{refs}

\end{document}
